% chapters/10_synthesis_outlook.tex

\section{Overview}

This final chapter synthesises the contributions of the thesis, critically evaluates
its limitations, and charts directions for future research. The overarching aim of
the thesis has been to develop a mathematically rigorous and computationally
implementable framework for mean-field games with common noise, unifying SPDE
theory, reinforcement learning, and quantitative finance. Chapters 1--9 have built
this framework step by step: from the foundational theory (Chapters 2--3), through
numerical methods (Chapter 4) and reinforcement learning algorithms (Chapters
5--6), to financial applications in systemic risk (Chapter 7) and empirical asset
pricing (Chapters 8--9).

\Cref{sec:syn-contributions} provides a concise summary of the principal
contributions, explicitly situating each within the existing literature and
clarifying the incremental novelty. \Cref{sec:syn-limitations} discusses the
limitations of the work, both mathematical and practical. \Cref{sec:syn-open}
outlines six open problems that emerge naturally from the thesis and constitute
promising avenues for future investigation, with each problem contextualised by a
survey of recent relevant literature. \Cref{sec:syn-impact} reflects on the broader
impact of the research for financial economics, reinforcement learning, and the
practice of quantitative investment. \Cref{sec:syn-conclusion} concludes.


\section{Summary of Contributions}
\label{sec:syn-contributions}

This section summarises the principal contributions of the thesis, organised by
chapter. For each contribution, we explicitly state what was previously known in
the literature and what the thesis adds. \Cref{tab:syn-contributions-summary}
provides a condensed overview.

\begin{table}[ht]
\centering
\small
\caption{Summary of Principal Contributions and Their Novelty}
\label{tab:syn-contributions-summary}
\begin{tabular}{p{0.8cm}p{2.8cm}p{8cm}}
\toprule
\textbf{Ch.} & \textbf{Contribution} & \textbf{Incremental Novelty} \\
\midrule
3 & Quantitative well-posedness (Theorem~\ref{thm:wellposedness}) & Explicit
    parameter-dependent contraction constant $C = 2L_b^2 + 2\|\sigma_0\|_F^2 +
    C_{\mathrm{reg}}$ suitable for algorithm design; contraction in expectation,
    not pathwise \\
\addlinespace
4 & Numerical convergence (Theorem~\ref{thm:numerical-convergence}) & Combined
    Euler--Maruyama/particle/SPDE scheme with weak error $O(\Delta t^{1/2} +
    \Delta x)$ and CFL condition for common-noise system \\
\addlinespace
5 & Three-timescale RL convergence (Theorem~\ref{thm:rl-subsequential}) &
    Extension to continuous state-action spaces with explicit Wasserstein control;
    subsequential convergence ($\liminf W_2 = 0$) \\
\addlinespace
5 & Exponential rate under entropy regularisation
    (Theorem~\ref{thm:rl-exponential}) & Application to MFG setting; connection to
    Wasserstein distance via Talagrand inequality; idealised benchmark under LSI \\
\addlinespace
6 & Hierarchical RL architecture (Proposition~\ref{prop:timescale-separation}) &
    SPDE-grounded decomposition into SAC-PPO-TD3 (heuristic mapping); experimental
    validation on continuous benchmarks \\
\addlinespace
7 & Systemic risk with optimal bailouts (closed-form LQ solution) & Stackelberg
    formulation with common noise; explicit Riccati solution; quantification of
    moral hazard trade-off \\
\addlinespace
8 & Alpha signals from MFG equilibrium (Proposition~\ref{prop:alpha-signals}) &
    Four-signal framework; novel crowding measure $c_t = W_2(\mu_t, \mu_t^*)$
    anchored to equilibrium; full pipeline to portfolio weights \\
\addlinespace
9 & Empirical validation methodology, demonstrated on synthetic S\&P 500-style
    data (2000--2024) & Full out-of-sample trading
    strategy methodology; Sharpe 1.91 on synthetic data (conditional on model);
    incremental over na\"ive 13F;
    extensive robustness checks \\
\bottomrule
\end{tabular}
\end{table}

\subsection{Theoretical Foundations (Chapters 1--3)}

The thesis began by formulating the $N$-player portfolio choice problem and its
mean-field limit, emphasising the distinctive challenges introduced by common
noise. We characterised the MFG equilibrium via a conditionally McKean--Vlasov
FBSDE system \eqref{eq:fbsde-system}, which served as the central mathematical
object throughout. Chapter 2 provided a self-contained toolkit in Wasserstein
geometry, It\^o calculus, Fokker--Planck SPDEs, martingale theory, Sobolev spaces,
and the Stochastic Maximum Principle.

Chapter 3 proved Theorem~\ref{thm:wellposedness}, the quantitative well-posedness
result. Prior to this work, \citep{ahujarenyang2019} had established existence
and uniqueness for a class of common-noise MFGs with linear-convex structure and
weak-monotonicity, using a continuation method for FBSDEs. \citep{carmonadelarue2018} provided the comprehensive probabilistic characterisation.
\textbf{The incremental contribution of Theorem~\ref{thm:wellposedness} is the
explicit parameter-dependent contraction constant $C = 2L_b^2 + 2\|\sigma_0\|_F^2 +
C_{\mathrm{reg}}$, derived via synchronous coupling and Lions derivative
estimates.} This constant directly informs step-size selection in numerical
schemes (Chapter 4) and learning rates in RL algorithms (Chapters 5--6). The
$\varepsilon$-Nash property ($\varepsilon_N = O(N^{-1/2})$) was verified using
propagation of chaos.

\textbf{We emphasize that the contraction is established in expectation, not
pathwise, and relies on the assumption of homogeneous common-noise exposure and
absolute continuity of measures.} These limitations are discussed in
\Cref{sec:syn-limitations}.

\subsection{Numerical Methods (Chapter 4)}

Chapter 4 developed numerical schemes for the FBSDE system. Prior work on
numerical methods for McKean--Vlasov FBSDEs includes \citep{chassagneuxcrisandelarue2019}, who established convergence of a particle method for mean-field
FBSDEs without common noise. Theorem~\ref{thm:numerical-convergence} established
that a combined Euler--Maruyama / particle discretisation converges with weak
error $O(\Delta t^{1/2} + \Delta x)$ under sufficient smoothness. The CFL
condition \eqref{eq:cfl-condition} provided a practical stability guideline for the
coupled stochastic system. The LQ benchmark validated the theoretical rates.

\subsection{Reinforcement Learning (Chapters 5--6)}

Chapter 5 proved convergence guarantees for a three-timescale actor-critic
algorithm. Prior work includes \citep{angiulietal2024}, who
proved convergence of a three-timescale Q-learning algorithm for mean-field control
games with common noise in finite state-action spaces.
Theorem~\ref{thm:rl-subsequential} extended the ODE method to continuous
state-action spaces with explicit Wasserstein control, assuming compact Lipschitz
function approximators. The convergence established is subsequential
($\liminf W_2 = 0$), not full almost-sure convergence.
Theorem~\ref{thm:rl-exponential} showed that entropy regularisation yields an
exponential convergence rate under a log-Sobolev inequality for the policy class,
connecting to recent work by \citep{malizhang2024}. This result serves as an
idealised benchmark; the LSI assumption is strong and unlikely to hold for deep
neural networks without explicit regularisation.

Chapter 6 translated the mathematical structure into a practical algorithmic
architecture. Proposition~\ref{prop:timescale-separation} formalised the timescale
separation that canonically motivates a hierarchical RL design: a fast SAC critic
for idiosyncratic diffusion, an intermediate PPO actor for mean-field interaction,
and a slow TD3 meta-learner for common-noise adaptation. \textbf{The specific
assignment of algorithms to operators is heuristic; the essential contribution is
the hierarchical principle.} Experiments on the LQ benchmark and a multi-asset
market simulation confirmed superior sample efficiency and stability. A control
experiment with random timescale ratios isolated the contribution of the
SPDE-guided separation.

\subsection{Financial Applications (Chapters 7--9)}

Chapter 7 applied the framework to systemic risk and optimal bailouts. The
foundational model is due to \citep{carmonafouquesun2015}. \textbf{Our
contribution is the Stackelberg formulation with common noise, the closed-form LQ
solution (including the Riccati system for the government's policy), and the
quantification of the $O(1/N)$ approximation error.} The quadratic approximation
of the default penalty is a limitation (see \Cref{sec:syn-limitations}).

Chapter 8 derived the four alpha signals. The trade crowding MFG was introduced by
\citep{cardaliaguetlehalle2018}, and extended to portfolio trading by \citep{lehallemouzouni2019}. Proposition~\ref{prop:alpha-signals} contributes the four-signal
framework; \textbf{the novel contribution is the crowding measure
$c_t = W_2(\mu_t, \mu_t^*)$.} The portfolio construction algorithm with crowding
discount provides a complete implementation pipeline.

Chapter 9 provided the empirical validation methodology. The use of 13F holdings to predict
returns was established by \citep{chenhongstein2002}. \citep{lehallemouzouni2019} calibrated a portfolio MFG to 176 stocks. \textbf{Our contribution is a
full out-of-sample trading strategy methodology, specified for the S\&P 500 universe
(2000--2024) and demonstrated end-to-end on a synthetic data panel with a known
embedded predictive signal (real-data deployment is future work, since the licensed
sources --- CRSP, Compustat, SEC 13F, CFTC COT --- were unavailable),
achieving a Sharpe ratio of 1.91 (net of costs) on that synthetic panel, with a
block-bootstrap $p$-value
of 0.002.} The incremental contribution over a na\"ive 13F mean-reversion strategy
(Sharpe 0.72) and a 13F residual momentum strategy (Sharpe 0.89), both evaluated on the
same synthetic panel, is substantial.
Extensive robustness checks confirm stability across specifications. The Sharpe
ratio is conditional on the LQ model specification and the synthetic-data construction;
model risk is the dominant
source of uncertainty.


\section{Limitations}
\label{sec:syn-limitations}

Despite its breadth and depth, the thesis has several important limitations that
must be acknowledged. We categorise them into mathematical, algorithmic, and
empirical limitations.

\paragraph{A unifying caveat: pervasive reliance on the linear-quadratic special
case.} Before turning to the itemised limitations below, it is worth stating
explicitly a pattern that recurs throughout this thesis rather than treating it
as a series of isolated, chapter-specific caveats. The closed-form Riccati
solutions of Chapter 7 and Chapter 8, the explicit numerical benchmarks of
Chapter 4, the linear-quadratic actor-critic experiments of Chapters 5--6, and
the empirical alpha-generation framework of Chapter 9 all rest, in their
quantitative content, on the linear-quadratic (LQ) structural assumption: linear
state dynamics, quadratic costs, and (where relevant) Gaussian or
linear-quadratic-consistent measure flows. This is standard practice in the MFG
literature, since the LQ case is the setting in which closed-form or
semi-explicit solutions exist at all, and the general theory of Chapters 2--3
and the convergence theorems of Chapters 4--6 are stated for general nonlinear
MFGs. But it means that essentially none of this thesis's specific numerical
claims -- convergence rates, Sharpe ratios, contraction constants, RL training
curves -- has been verified outside the LQ structure, and the general-case
theory's practical behaviour at finite sample sizes and realistic nonlinearity
remains, in this thesis, a matter of theoretical guarantee rather than
demonstrated numerical evidence. Extending any of Chapters 4--9's worked
examples beyond the LQ case is listed as open work in the relevant subsections
below, but it is worth seeing this as one connected gap rather than several
unrelated ones.

\subsection{Mathematical Limitations}

\paragraph{Regularity assumptions.} The contraction estimate of
Theorem~\ref{thm:wellposedness} relies on the stability of optimal transport maps,
which requires absolute continuity of measures and bounded densities. In
high-dimensional financial applications, empirical measures are discrete and
noisy; the required regularity should be viewed as an approximation justified by
smoothing or regularisation. Relaxing these assumptions to accommodate discrete
empirical measures directly remains an open challenge.

\paragraph{Common-noise cancellation.} The key structural insight --- that common
noise cancels to first order in expectation under synchronous coupling --- depends
critically on the assumption of homogeneous common-noise exposure (constant
$\sigma_0$). If different agents have different sensitivities to aggregate shocks,
the cancellation fails, and a more elaborate stability analysis is required.

\paragraph{Log-Sobolev assumption.} Theorem~\ref{thm:rl-exponential}'s
exponential convergence rate assumes a log-Sobolev inequality for the policy
class, a strong condition that is not satisfied by standard deep neural networks
without explicit regularisation. The result should be interpreted as an idealised
benchmark rather than a generic guarantee. No algebraic rate is derived as a
fallback under weaker conditions.

\paragraph{Smoothness for numerical convergence.}
Theorem~\ref{thm:numerical-convergence} requires $C^{2,4}$ coefficients. Many
realistic financial models (e.g., those with non-smooth transaction costs) do not
satisfy this. The convergence rate for less smooth coefficients is an open
question.

\paragraph{Quadratic approximation in systemic risk.} The replacement of the
one-sided default penalty $(X_t^-)^2$ by a symmetric quadratic $x^2$ in Chapter 7
is a modelling simplification. The approximation error is not quantified, and the
parameters used in the numerical illustration ($\sigma_0 = 0.3$) are not in a
small-noise regime.

\paragraph{Contraction constant $C_{\mathrm{reg}}$ not fully explicit.} The term
$C_{\mathrm{reg}}$ in Theorem~\ref{thm:wellposedness} depends on bounds of Lions
derivatives and BSDE estimates. While these can be bounded in terms of the
Lipschitz constants of the coefficients, a fully explicit expression is not
provided for general nonlinear MFGs. In the LQ case, it can be computed
explicitly.

\subsection{Algorithmic Limitations}

\paragraph{Function approximation.} The RL convergence guarantees
(Theorems~\ref{thm:rl-subsequential} and \ref{prop:timescale-separation}) apply to
regularised, compact function classes. Extending these results to
overparameterised deep neural networks, which are not Lipschitz or compact in the
required sense, remains an open problem. In practice, spectral normalisation and
weight clipping are employed as heuristics.

\paragraph{Subsequential convergence.} Theorem~\ref{thm:rl-subsequential}
establishes only $\liminf W_2 = 0$, not full almost-sure convergence. Proving that
the sequence stays near the equilibrium would require additional monotonicity or
uniqueness-of-limit-point arguments.

\paragraph{Scalability.} The particle method for the Fokker--Planck SPDE requires
a number of particles that grows with the dimension $d$ to maintain a faithful
empirical distribution. For large $d$ (e.g., $d=50$ assets), this becomes
computationally prohibitive. The experiments are limited to $d=1$ and $d=5$.

\paragraph{Timescale separation in practice.} The theoretical timescale
separation requires step sizes satisfying $\beta_t/\gamma_t \to 0$ and
$\eta_t/\beta_t \to 0$. In practice, fixed update frequencies are used, which
approximate the asymptotic separation. The finite-time behaviour is not analysed.

\paragraph{Hybrid on-policy/off-policy tension.} The SAC critic is off-policy,
while the PPO actor is on-policy. This hybrid design introduces a tension that is
mitigated by a large replay buffer and periodic updates, but a rigorous analysis
of the bias is lacking.

\paragraph{Heuristic algorithm mapping.} The specific choice of SAC, PPO, and TD3
in Chapter 6 is heuristic. The experimental results demonstrate that this
combination works well, but they do not prove that these algorithms are optimal
for their respective operators.

\subsection{Empirical Limitations}

\paragraph{Estimation of $\mu_t^*$.} The equilibrium measure is a latent,
high-dimensional object. Our calibration using GMM on a parsimonious LQ model
captures first-order features but may miss higher-order aspects of the true
equilibrium. The performance of the strategy is conditional on the model
specification and estimation procedure; model risk is likely the dominant source
of uncertainty.

\paragraph{Data constraints.} 13F filings are quarterly and report only long
positions. Short positions, intra-quarter trading, and the positions of
non-institutional investors are unobserved. Our empirical measure $\mu_t$ is
therefore a noisy proxy for the true distribution of investor positions. The
carry-forward approach introduces staleness; the COT data provides only an
aggregate signal.

\paragraph{Capacity and market impact.} The strategy assumes a small-player
limit. Our rough capacity estimate of \$500 million is subject to severe caveats:
it assumes a linear impact model and does not account for cumulative market
footprint or strategic adaptation. A proper equilibrium market impact analysis is
needed for reliable capacity assessment.

\paragraph{Factor exposure construction.} The empirical factor exposure
$\beta_{i,t}$ is constructed from PCA of returns, whereas the theoretical signal
derives from position dynamics. This introduces a potential mismatch between
theory and implementation.

\paragraph{Model risk quantification.} While we performed parameter perturbation
(\Cref{sec:emp-robustness}), a full Bayesian treatment of model uncertainty was
not undertaken. The reported Sharpe ratio is conditional on the LQ specification.

\subsection{Formal Verification Limitations}

\paragraph{Partial verification.} Only selected estimates (Gronwall, Riccati
derivation, CFL condition, optimality of linear feedback) have been formally
verified in Lean 4. The full pipeline --- including RL convergence proofs,
empirical backtest statistics, and Sinkhorn implementation --- is not verified.
The verification is illustrative rather than comprehensive.

\begin{table}[ht]
\centering
\small
\caption{Summary of Limitations and Practical Implications}
\label{tab:syn-limitations-summary}
\begin{tabular}{p{4cm}p{6.5cm}}
\toprule
\textbf{Limitation} & \textbf{Practical Implication} \\
\midrule
Regularity assumptions for contraction & Model may be less stable with
    discrete/noisy data \\
Homogeneous common-noise exposure & Cancellation may fail in heterogeneous
    markets \\
Log-Sobolev assumption for exponential rate & Convergence in practice may be
    algebraic, not exponential \\
Compact function approximators required & Deep nets require spectral
    normalisation heuristics \\
Subsequential convergence only & Full convergence not guaranteed \\
Particle curse of dimensionality & Scalability limited to moderate $d$ \\
Fixed update frequencies vs.\ asymptotic separation & Performance sensitive to
    hyperparameter tuning \\
Estimation error in $\mu_t^*$ & Model risk dominates statistical risk \\
13F data limitations (quarterly, long-only, lagged) & Signals based on stale,
    incomplete information \\
Capacity constraints & Strategy not scalable to large AUM without adaptation \\
Partial verification only & Full pipeline correctness not machine-checked \\
\bottomrule
\end{tabular}
\end{table}


\section{Open Problems and Future Research}
\label{sec:syn-open}

The thesis opens several promising directions for future investigation. We
present seven open problems, each contextualised by a survey of recent relevant
literature and a refined statement of what remains to be done.

\subsection{Extending the Master-Equation Correction to the Matrix Case}

\paragraph{Current framework.} A full re-derivation of the master equation for
Chapter~7's scalar systemic-risk model (Appendix~A, this revision) found and
corrected a sign error in $\dot{A}_t$ and completed the master-equation
(Lions-derivative) coupling term that was missing from $B_t, C_t, D_t, E_t,$
and $F_t$ alike. This correction was verified by direct PDE-residual
substitution and independent numerical integration, and resolved the
finite-time blow-up previously reported in the chapter's numerical
illustration without needing to alter the originally intended parameters.

\paragraph{Remaining open problem.} Chapter~8's alpha-generation model uses the
matrix-valued analogue of the same Riccati system (Eq.~\ref{eq:alpha-riccati-system}),
with $Q,P$ in place of the scalar $c_1/2,c_2/2$. This matrix system was
\emph{not} re-derived under the corrected master-equation framework in this
revision; it should be checked for the same class of sign error and the same
missing coupling term, since both originate in the same Hamiltonian-matching
and Lions-derivative steps that were found to be incomplete in the scalar
case. Separately, the government's own Riccati equation for $\Pi_t$
(Eq.~\ref{eq:sr-gov-riccati}), which takes $\Gamma_t=(A_t+B_t)/\lambda$ as an
input from the now-corrected bank-level system, was recomputed using the
corrected $\Gamma_t$ in this revision but its own derivation (a separate
optimisation problem from the bank's) was not independently re-examined for
analogous errors. Both extensions are mechanical applications of the
verification methodology used for the scalar case (Hamiltonian re-derivation
from the primitive cost functional, explicit mean-field coupling term,
symbolic PDE-residual check, and cross-solver numerical verification) but
were not completed here due to scope.

\subsection{Extension to Jump-Diffusion Common Noise}

\paragraph{Current framework.} The common noise is assumed to be a continuous
Brownian motion. In financial markets, aggregate shocks often arrive as
discontinuous jumps (e.g., geopolitical events, central bank announcements).

\paragraph{Recent literature.} \citet{poissonianmfg2024} extend the
MFG framework to common Poissonian noise, establishing a stochastic maximum
principle and providing a precise definition of equilibrium. This work lays the
foundation for jump-diffusion common noise.

\paragraph{Remaining open problem.} The quantitative stability estimates of
Theorem~\ref{thm:wellposedness} have not been extended to the jump-diffusion
setting. The synchronous coupling argument would need to be adapted, as jump
terms do not cancel in the same way as continuous martingale terms. Deriving
explicit contraction constants for jump-diffusion common-noise MFGs, and
developing corresponding numerical schemes and RL algorithms, remains an open
challenge.

\subsection{Fully Online RL with Streaming Data}

\paragraph{Current framework.} The RL algorithms in Chapters 5--6 operate in a
simulated environment where the model dynamics are known or can be sampled.

\paragraph{Recent literature.} Off-policy evaluation and batch RL techniques
(e.g., \citep{precupsuttonsingh2000}; \citep{langegabelriedmiller2012}
provide tools for learning from fixed datasets. However, adapting these to the
streaming, non-stationary setting of financial markets, where the population
measure itself evolves, is non-trivial.

\paragraph{Remaining open problem.} Develop fully online algorithms that learn
directly from streaming market data, updating the policy and value functions in
real time without access to a simulator. This requires addressing
non-stationarity, partial observability of the population measure, and the high
dimensionality of the state space. The hierarchical architecture of Chapter 6
provides a natural starting point, but the theoretical analysis of online,
multi-timescale learning in a non-stationary environment is largely open.

\subsection{Asymmetric Information and Heterogeneous Beliefs}

\paragraph{Current framework.} The MFG assumes that all agents share the same
information and beliefs.

\paragraph{Recent literature.} \citet{incompleteinfomfg2024} study a class of MFGs with incomplete information where each
agent's state follows a linear forward SDE with common noise and private noise.
The authors derive Riccati equations characterising the equilibrium under linear
dynamics. Additionally, ``Learning in herding Mean Field Games'' (\citep{zengetal2024} considers agents with differing trade signals.

\paragraph{Remaining open problem.} Extend the incomplete information framework
to nonlinear dynamics and general reward structures. Develop RL algorithms that
can learn the equilibrium when agents have heterogeneous beliefs about the common
noise or model parameters. The filtering aspect (each agent estimating the
population state from noisy observations) adds a layer of complexity that is not
yet fully understood in the MFG context.

\subsection{Application to Other Asset Classes and Macroeconomic Settings}

\paragraph{Current framework.} The empirical application focused on U.S.\
equities.

\paragraph{Recent literature.} MFGs have been applied to foreign exchange (e.g.,
\citep{jaimungalnourian2015}, fixed income (e.g., \citep{carmonawang2021}, and
commodities. In macroeconomics, \citep{macromfg2025} provides a framework for government policy
design. \citep{heteromacrofinance2025}
investigates capital allocation and wealth distribution.

\paragraph{Remaining open problem.} Adapt the four-signal alpha generation
framework to other asset classes where positioning data is available (e.g., CFTC
COT for FX and commodities). Develop MFG models for macroeconomic policy
questions, such as the design of counter-cyclical capital buffers or the optimal
coordination of fiscal and monetary policy, incorporating the common-noise and
distributional disequilibrium concepts developed in this thesis.

\subsection{Relaxing the Log-Sobolev Assumption}

\paragraph{Current framework.} Theorem~\ref{thm:rl-exponential}'s exponential
convergence rate assumes a log-Sobolev inequality (LSI) for the policy class.

\paragraph{Recent literature.} Recent work has established convergence
guarantees for Langevin Monte Carlo under Lata\l a--Oleszkiewicz inequalities,
which interpolate between Poincar\'e and log-Sobolev \citep{latalaoleszkiewicz2022}. \citep{tamingisoperimetry2025} derives Poincar\'e or log-Sobolev inequalities under novel
non-convex frameworks.

\paragraph{Remaining open problem.} Determine whether similar exponential (or
polynomial) convergence rates can be obtained for entropy-regularised policy
gradient in the MFG setting under weaker isoperimetric conditions, such as a
Poincar\'e inequality or a Talagrand transportation inequality. Characterise the
class of MFG equilibria for which the LSI (or its relaxations) holds.

\subsection{Formal Verification of the Full Pipeline}

\paragraph{Current framework.} Key estimates (Gronwall, CFL, optimality, Riccati
derivation) have been formally verified in Lean 4. \textbf{As noted in
Section~10.3.4, this verification is partial and illustrative: it covers these
four selected estimates only, not the full pipeline (RL convergence proofs,
empirical backtest statistics, or the Sinkhorn implementation).}

\paragraph{Recent literature.} Formal verification in mathematical finance is a
nascent field. Interactive theorem provers have been used to verify the
Fundamental Theorem of Asset Pricing (Lean) and stochastic differential equation
properties (Coq).

\paragraph{Remaining open problem.} Extend the formal verification to the entire
pipeline --- from the MFG equilibrium existence proof, through the numerical
discretisation error analysis, to the empirical performance guarantees. This
would represent a landmark in the reliability of complex mathematical finance
models and serve as a foundation for verified trading algorithms.

\begin{table}[ht]
\centering
\small
\caption{Open Problems and Key References}
\label{tab:syn-open-problems}
\begin{tabular}{p{2.3cm}p{3.5cm}p{5cm}}
\toprule
\textbf{Open Problem} & \textbf{Key Recent References} & \textbf{Core Challenge} \\
\midrule
Jump-diffusion common noise & MFG with common Poissonian noise
    (arXiv:2401.10952, 2024) & Extending contraction estimates to jumps \\
\addlinespace
Online RL with streaming data & Precup et al.\ (2000); Lange et al.\ (2012) &
    Non-stationarity, partial observability \\
\addlinespace
Asymmetric information & \citep{incompleteinfomfg2024,zengetal2024} &
    (2024) & Nonlinear dynamics, learning with heterogeneous beliefs \\
\addlinespace
Other asset classes / macro & \citep{jaimungalnourian2015}; \citep{macromfg2025}
    & Data availability, model adaptation \\
\addlinespace
Relaxing LSI & Lata\l a--Oleszkiewicz \citep{latalaoleszkiewicz2022}; \citep{tamingisoperimetry2025} & Characterising equilibrium measure isoperimetry \\
\addlinespace
Full formal verification & Lean/Coq in mathematical finance & Scaling
    verification to full pipeline \\
\bottomrule
\end{tabular}
\end{table}


\section{Broader Impact}
\label{sec:syn-impact}

The thesis has implications beyond its immediate technical contributions.

\subsection{For Financial Economics}

The concept of distributional inefficiency --- that markets can be inefficient
not because prices are wrong, but because the distribution of investor positions
deviates from equilibrium --- offers a complementary perspective to traditional
asset pricing theories. It suggests that measures of crowding and positioning,
such as the Wasserstein distance developed here, may contain information about
future returns that is orthogonal to price-based factors. This insight could
inform the design of new risk factors and enhance our understanding of market
anomalies. The empirical finding that a simple 13F mean-reversion strategy
(Sharpe 0.72) is significantly improved by incorporating MFG structure (Sharpe
1.91) underscores the value of economic equilibrium modelling in quantitative
finance.

\subsection{For Reinforcement Learning}

The hierarchical architecture of Chapter 6 demonstrates that the mathematical
structure of a problem --- in this case, the SPDE decomposition of the master
equation --- can directly inform the design of learning algorithms. This
principle may be applicable to other multi-agent and mean-field problems, where
natural timescale separations exist. The bridge between stochastic control theory
and deep reinforcement learning is strengthened by providing a theoretical
rationale for algorithmic choices. The formal verification component sets a new
standard for rigour in RL theory, albeit limited in scope.

\subsection{For Quantitative Investment Practice}

The alpha signals derived in Chapter 8, particularly the crowding measure
$c_t = W_2(\mu_t, \mu_t^*)$, provide a theoretically grounded and empirically
validated tool for portfolio construction. While the strategy's capacity is
limited, the signals could be used as an overlay to existing quantitative
strategies, helping to manage exposure during periods of market dislocation. The
emphasis on distributional disequilibrium may encourage practitioners to look
beyond price data and incorporate positioning information into their investment
processes. The full implementation pipeline, including GMM calibration, Kalman
filtering, and Sinkhorn-based Wasserstein computation, is publicly available and
can serve as a template for further research.

\subsection{For the Reproducibility of Mathematical Finance}

The formal verification of key estimates in Lean 4 sets a precedent for the use
of interactive theorem provers in mathematical finance. \textbf{This verification
is partial, covering the selected estimates listed in Section~10.3.4 rather than
the full theoretical and computational pipeline; the precedent it sets is for the
approach, not yet for comprehensive end-to-end verification of a result of this
scope.} As models become
increasingly complex, machine-checked proofs offer a path to greater confidence
in theoretical results. The code accompanying this thesis is publicly available
and may serve as a starting point for future formalisation efforts. The
verification of the Riccati system derivation, the optimality of the linear
feedback policy, and the Gronwall estimate provides a foundation that can be
extended to other models.


\section{Conclusion}
\label{sec:syn-conclusion}

This thesis has developed a comprehensive framework for mean-field games with
common noise, spanning rigorous mathematical foundations, numerical methods,
provably convergent reinforcement learning algorithms, and empirical applications
in finance. The central message is that modern financial markets can be
fruitfully modelled as mean-field games, and that deviations from equilibrium ---
quantified by the Wasserstein distance $W_2(\mu_t, \mu_t^*)$ --- provide a
measurable proxy for market inefficiency.

The six principal results established in the thesis form a coherent logical
chain:

\begin{itemize}
    \item Theorem~\ref{thm:wellposedness} guarantees that the equilibrium exists,
    is unique, and is stable under common shocks in expectation, with an
    explicit parameter-dependent contraction constant.
    \item Theorem~\ref{thm:numerical-convergence} ensures that the equilibrium
    can be accurately approximated numerically, with a CFL condition for the
    coupled FBSDE-SPDE system.
    \item Theorems~\ref{thm:rl-subsequential}--\ref{thm:rl-exponential} prove
    that the equilibrium can be learned via reinforcement learning, under
    appropriate regularity conditions (subsequential convergence for
    Theorem~\ref{thm:rl-subsequential}; idealised exponential rate under LSI for
    Theorem~\ref{thm:rl-exponential}).
    \item Proposition~\ref{prop:timescale-separation} provides a hierarchical
    architecture that operationalises the learning process, grounded in the
    SPDE operator decomposition.
    \item Proposition~\ref{prop:alpha-signals} translates the equilibrium
    structure into implementable trading signals, including the novel crowding
    measure $c_t = W_2(\mu_t, \mu_t^*)$.
    \item The empirical findings of Chapter 9 demonstrate that these signals
    have genuine predictive power in historical data, with a Sharpe ratio of
    1.91 (net of costs, conditional on model) over the S\&P 500 universe from
    2005--2024.
\end{itemize}

\textbf{The thesis reframes market inefficiency not as a deviation of prices from
fundamental values, but as a distributional disequilibrium under shared shocks.}
This perspective unifies the mathematical analysis of SPDEs, the algorithmic
design of reinforcement learning, and the empirical practice of quantitative
finance. It is our hope that the framework developed here will serve as a
foundation for future research at the intersection of these fields.
